% Philippica IV (philippic-4) --- English translation
% Translated by Claude (Anthropic) from the Latin (Perseus).
% Style guide: STYLE.md.

\ciceroSection{1} Your incredible crowd, citizens, and an assembly as great as I do not seem to remember, both brings me the greatest eagerness to defend the commonwealth and the greatest hope of recovering it. Spirit, indeed, has never failed me; what failed were the times --- and as soon as they seemed to show some glimpse of light, I was first to defend your liberty. If I had tried to do it earlier, I could not be doing it now. For on this day, citizens, do not suppose that some trivial matter has been done: the foundations have been laid for the actions still to come. For Antony has not yet been called an enemy by the Senate in word, but in fact he has already been so judged.

\ciceroSection{2} Now indeed I am far more lifted up, since you too have approved, with so great a consent and so great an outcry, that he is an enemy. For it cannot be, citizens, that those men are not impious who have raised armies against the consul, or else that he is not an enemy against whom arms have rightly been taken up. This doubt, then, although there was none, the Senate on this day removed, so that none could remain. Gaius Caesar, who has defended and is defending the commonwealth and your liberty by his own zeal, his own counsel, and finally his own patrimony, has been honoured with the highest praises of the Senate.

\ciceroSection{3} I praise you, I praise you, citizens, that with the most grateful hearts you attend the name of that most illustrious young man --- or boy, rather; for his deeds are those of immortality, his name those of his age. Many things I remember, many I have heard, many I have read, citizens: nothing of this kind have I learned of from the memory of all the ages: that, when we were pressed down by slavery, when the evil was growing day by day, when we had no protection, when we feared the deadly and pestilent return of Marcus Antonius from Brundisium, this young man should have taken a plan unhoped-for by all and certainly unheard of --- to raise an unconquered army out of his father's soldiers, and to turn aside from the ruin of the commonwealth the frenzy of Antony, goaded on by the cruellest counsels.

\ciceroSection{4} For who is there who does not understand this: that, had Caesar not prepared an army, Antony's return would not have been without our destruction? For he was returning so ablaze with hatred against you, so steeped in the blood of the Roman citizens whom he had slaughtered at Suessa, whom at Brundisium, that he was thinking of nothing but the ruin of the Roman people. What protection, then, was there for your safety and your liberty, had there not been the army of Gaius Caesar, raised from the bravest of his father's soldiers? Concerning his praises and the honours which, in return for his more-than-divine and immortal services, are owed to him as divine and immortal, the Senate with my assent decreed a little while ago that the matter be brought forward at the earliest opportunity.

\ciceroSection{5} Who, by that decree, does not see that Antony has been judged an enemy? For what else can we call him, when the Senate judges that singular honours must be sought out for those who lead armies against him? What, again, of the Martian Legion --- which seems to me to have drawn its name, by divine appointment, from that god from whom we are taught the Roman people was begotten --- did it not, by its own decrees, judge Antony an enemy before the Senate did? For if he is no enemy, we must of necessity judge these men enemies who have abandoned the consul. Splendidly and aptly, citizens, with your shouts of acclamation you have ratified the most glorious deed of the Martian soldiers: they who attached themselves to the authority of the Senate, to your liberty, to the whole commonwealth, abandoned that man, an enemy, a brigand, and a parricide of his country.

\ciceroSection{6} Nor did they do this with spirit and bravery only, but also with deliberation and wisdom: they took their stand at Alba --- a town well placed, fortified, near at hand, of the bravest of men, of the most loyal and best of citizens. Imitating the courage of this legion, the Fourth Legion, under the lead of Lucius Egnatuleius --- whom the Senate, deservedly, has just praised --- has followed the army of Gaius Caesar. What graver judgments are you waiting for, Marcus Antonius? Caesar is exalted to the skies for raising an army against you; the legions which abandoned you, which had been sent for by you, which would have been yours had you preferred to be a consul rather than an enemy, are praised in the choicest words; the bravest and truest judgment of those legions the Senate confirms, the whole Roman people approves --- unless perhaps you, citizens, judge Antony a consul, not an enemy.

\ciceroSection{7} So I thought, citizens, that you would judge, as indeed you show. What, then? Do you think that the towns, the colonies, the prefectures judge any differently? All mortals consent with one mind; all arms are to be taken up by those who would have these things safe, against that plague. What, again, of the judgment of Decimus Brutus, citizens --- which you have been able to perceive from his edict of this very day --- does it seem to anyone, after all, fit to be despised? Rightly and truly do you deny it, citizens. For it is as though, by the kindness and gift of the immortal gods, the race and name of the Bruti has been granted to the commonwealth for either the establishing or the recovering of the liberty of the Roman people.

\ciceroSection{8} What, then, has Decimus Brutus judged concerning Marcus Antonius? He shuts him out of his province; he resists him with an army; he urges all Gaul to war --- Gaul itself roused of its own accord and by its own judgment. If Antony is consul, Brutus is an enemy: if Brutus is the preserver of the commonwealth, Antony is an enemy. Can we, then, doubt which of these is the case? And just as you, with one mind and one voice, deny that you doubt, so the Senate has just now decreed that Decimus Brutus deserves the best of the commonwealth, in defending the authority of the Senate and the liberty and empire of the Roman people. From whom was he defending them? From an enemy, of course: for what other defence is there that deserves praise?

\ciceroSection{9} Next the province of Gaul is praised, and rightly adorned by the Senate with the fullest words, because it resists Antony. If that province thought him consul and would not receive him, it would bind itself by a great crime: for all the provinces ought to be in the law and power of the consul. Decimus Brutus, commander, consul-elect, a citizen born for the commonwealth, denies this; Gaul denies it; all Italy denies it; the Senate denies it; you deny it. Who, then, takes him for a consul, except the brigands? --- though not even they, in fact, feel what they say, nor can they dissent from the judgment of all mortals, however impious and wicked they are (and they are). But the hope of plundering and seizing blinds the minds of those whom no donation of goods, no assignment of lands, no spear of those endless auctions has ever satisfied; who have set out the City, the goods and fortunes of the citizens, as their plunder; who suppose, so long as there is anything here for them to seize and carry off, that they shall lack for nothing --- to whom Marcus Antonius (O immortal gods, avert and curse this omen, I beseech you!) has promised that he will divide the City among them.

\ciceroSection{10} So indeed, citizens, may it fall out for them as you pray, and may the penalty for this madness recoil upon him and his household! That so it shall be, I am confident. For now I judge that not only men but the immortal gods themselves have agreed to preserve the commonwealth. For whether by prodigies and portents the immortal gods foretell the future to us, things have been so openly proclaimed that penalty draws near to him and liberty to us; or whether so great a consent of all could not have been without the prompting of the gods, what is there that we could doubt about the will of the heavens?

\ciceroSection{11} It remains, citizens, for you to persevere in that resolve which you carry openly before you. I shall therefore do as commanders, with the battle-line drawn up, are accustomed to do: although they see their soldiers most ready for the fight, they still exhort them; so I shall encourage you, burning and lifted up for the recovery of liberty. Yours is no contest, citizens, with an enemy with whom any terms of peace can be possible. He no longer covets your slavery, as before, but, now enraged, your blood. No sport seems to him more delightful than gore, than slaughter, than the butchery of citizens before his eyes.

\ciceroSection{12} You have to do, citizens, not with a wicked and impious man, but with a vast and foul beast which, since it has fallen into the pit, must be buried under it. For if he climbs out of there, no cruelty of punishment will be too much to refuse. But he is held, he is hemmed in, he is pressed --- now by the forces we have already, soon by those which the new consuls will assemble in a few days. Lean upon the cause, citizens, as you are doing. Never was your consent greater in any cause; never were you so vehemently allied with the Senate. And no wonder: for it is not on what terms we shall live that is at stake, but whether we shall live at all, or perish with torment and disgrace.

\ciceroSection{13} Death, indeed, nature has set before all; but the cruelty and dishonour of death is for valour to ward off --- valour which is the property of the Roman race and stock. Hold fast to this, I beg, citizens, which your ancestors left you as an inheritance. All other things are false, uncertain, perishable, shifting: valour alone is set with the deepest roots; she can never by any force be shaken, never moved from her place. By her your ancestors first conquered all Italy, then razed Carthage, overthrew Numantia, brought into the power of this empire the most powerful kings and the most warlike peoples.

\ciceroSection{14} And your ancestors, citizens, had to do with an enemy who had a commonwealth, a senate-house, a treasury, the consent and concord of citizens --- with whom, if circumstances so required, there was at least some prospect of peace and treaty. This enemy of yours assails your commonwealth, while himself having none; he longs to destroy the Senate (that is, the council of the whole world), while himself having no public council; he has emptied your treasury, while having none of his own. For how can a man have the concord of citizens, when he has no citizen body at all? And what prospect of peace can there be with one in whom there is incredible cruelty and no good faith?

\ciceroSection{15} Yours, then, citizens --- the Roman people, conqueror of all nations --- the whole contest is with an assassin, with a brigand, with a Spartacus. For as for his being accustomed to glory that he is like Catiline, he is equal to him in crime, inferior in energy. Catiline, having no army, suddenly threw one together: this man has lost the army which he received. As, then, you broke Catiline by my diligence, by the Senate's authority, by your own zeal and valour, so the wicked banditry of Antony you shall in a short time hear has been crushed by your concord with the Senate (greater than ever there was), by the good fortune and valour of your armies and their commanders.

\ciceroSection{16} For my part, so much as I can accomplish and achieve by care, by labour, by sleepless nights, by authority, by counsel, I shall pass over nothing that I think pertains to your liberty: for indeed, in view of your most generous services to me, I cannot do otherwise without crime. This day, with the motion brought by a most brave man and a friend most devoted to you, this Marcus Servilius, and by his colleagues, men most distinguished, citizens most excellent, for the first time after a long interval, with me as proposer and leader, we have caught fire with the hope of liberty.
